\section{Methodology}
\label{sec:method}

\subsection{Datasets}

\paragraph{DAiSEE.}
The Dataset for Affective States in E-Environments
(DAiSEE)~\cite{gupta2016daisee} contains 9{,}068 ten-second video clips
from 112 subjects recorded via webcam during e-learning sessions.  Each
clip is annotated on a four-point Likert scale for engagement (0 = not
engaged, 3 = highly engaged) along with boredom, confusion, and
frustration.  We focus exclusively on engagement.  The official test
split comprises 1{,}784 clips across 21 subjects; the label distribution
is heavily skewed, with only 4 samples at level 0 and approximately 850
each at levels 2 and 3.

We evaluate on a stratified sample of $N\!=\!300$ test clips, retaining
all 4 level-0 clips, all 84 level-1 clips, and randomly sampling 106
clips each from levels 2 and 3.  From each 10-second clip, we extract
a single representative frame at $t\!=\!5$\,s (the temporal midpoint)
using \texttt{ffmpeg}.  This setup tests whether static visual cues
visible in a single webcam frame are sufficient for zero-shot engagement
inference.

\paragraph{Student Classroom Behaviour (SCB).}
The SCB dataset~\cite{wang2023scb} contains classroom scene images with
per-student bounding box annotations across six behaviour categories:
\emph{reading}, \emph{writing}, \emph{raising hand}, \emph{using phone},
\emph{bowing head}, and \emph{leaning over table}.  We use a combined
split of 1{,}168 validation images drawn from both the original 5K
corpus and an additional university subset.

Following the aggregation strategy proposed in the dataset documentation,
we derive a scene-level engagement label from the per-student annotations.
Each behaviour is mapped to an engagement score (0--3): disruptive
behaviours (\emph{phone}, \emph{bowing head}) score 0; passive
disengagement (\emph{leaning over}) scores 1; productive behaviours
(\emph{reading}, \emph{writing}) score 2; and active participation
(\emph{raising hand}) scores 3.  The scene-level label is the
mean across all annotated students, discretised via uniform thresholds
($\left[0,\,0.75\right) \to 0$; $\left[0.75,\,1.75\right) \to 1$;
$\left[1.75,\,2.75\right) \to 2$; $\left[2.75,\,3\right] \to 3$).
The resulting label distribution is highly imbalanced: 54.5\% of images
fall at level 3, 33.6\% at level 2, 10.1\% at level 1, and 1.9\%
at level 0.

\subsection{Models}

We evaluate four models at substantially different points on the
capability-cost spectrum:

\textbf{CLIP (ViT-B/32)}~\cite{radford2021learning} is a discriminative
contrastive model.  Classification proceeds by computing cosine
similarity between the image embedding and embeddings of textual
descriptions of each engagement level; the class with the highest
similarity is the prediction.  Inference is fully deterministic.

\textbf{BLIP-VQA (blip-vqa-base)}~\cite{li2022blip} is a discriminative
VQA model.  We pose a natural-language question about student engagement;
the model generates a short free-text answer (\eg, ``not very'',
``yes'') that is mapped to a level via a hand-crafted lookup table.
Beam decoding with $k\!=\!3$ ensures deterministic output.

\textbf{GPT-4o}~\cite{openai2024gpt4o} is a proprietary closed-source
multimodal model accessed via the OpenAI API.  Images are submitted as
base64-encoded JPEG at ``low'' detail resolution.  Temperature is 0.0
for run 1 and 0.7 for subsequent consistency runs.

\textbf{LLaVA-1.5-7B}~\cite{liu2023visual} is an open-source instruction-tuned
model served locally through Ollama at 4-bit quantisation (reducing
memory from ${\sim}14$\,GB to ${\sim}4$\,GB) on a local machine with 16\,GB RAM.
We use the same prompts as GPT-4o to enable direct comparison.
Temperature follows the same protocol as GPT-4o.

\subsection{Prompt Variants}

Three prompts are applied to each model to quantify sensitivity:

\textbf{P1 (Minimal)} provides only a four-label numeric scale with no
contextual guidance.  It establishes a lower bound on what each model
can infer from minimal language grounding.

\textbf{P2 (Rubric-Anchored)} augments P1 with explicit behavioural
anchors for each level, \eg, \emph{``looking away, yawning''} for
level 0 and \emph{``leaning forward, showing curiosity''} for level 3.
This mirrors how human annotators are typically trained.

\textbf{P3 (Chain-of-Thought)} instructs the model to first describe
the student's or class's observed facial expression and body posture
before committing to a rating.  This structured reasoning format has
shown benefits in language-only tasks~\cite{wei2022chain} and is
evaluated here in the visual domain.

For CLIP, the prompts are class-specific text templates.  For BLIP-VQA,
the prompts are question strings.  For GPT-4o and LLaVA, identical
prompt text is used to enable direct cross-model comparison.

\subsection{Evaluation Metrics}

All experiments are evaluated with four metrics: (i)~\textbf{Accuracy}
(exact-match 4-class); (ii)~\textbf{F1 Macro}, the unweighted mean of
per-class F1 scores, which penalises class collapse; (iii)~\textbf{Cohen's
$\kappa$ (quadratic weighted)}, which accounts for ordinal agreement
structure and penalises predictions far from ground truth more severely;
and (iv)~\textbf{MSE} (mean squared error), measuring average ordinal
distance.  For self-consistency, we report full 3-run agreement rate
and pairwise agreement rate.
